\section{Conclusion}

Exchanging only scalar scores can make distributed ES repeat substantial work.
Splitting update reconstruction helps when the computation saved exceeds the cost of
communicating a model-sized update. In our single-host measurements, this favors
splitting for full-rank perturbations and replication for low-rank perturbations at
large matrix width. The same preference holds on Qwen2.5-0.5B on the TPU. Across the
measured slow connection between A100 hosts, large full-rank updates favor replication too.

To choose a placement, measure reconstruction and the added communication cost in the
relevant computation context. Isolated communication benchmarks and ideal $1/D$ compute
scaling can misestimate the small differences that decide low-rank cases. The model
explains most measured preferences but is not an exact decision rule.

Low-rank perturbations also reduce generation time relative to stored full-rank noise
on the block, especially on the tested TPU. Their updates have similar measured
directional alignment, and the tested ranks reach similar observed final Countdown
accuracy. Full rank learns faster per sample at the first checkpoint. The limited
learning experiments support using the cheaper perturbations in this setting, while
leaving broader learning-efficiency comparisons open.
